\section*{Chapitre \ref{ch:rc:rates-risk} --- Risque de taux}

\begin{solution}{exo:rc:rates-risk:1}
DV01 parallèle $50-120+80 = +10$ milliers de dollars par point de base. P\&L
$50\times(-5) - 120\times2 + 80\times6 = -250-240+480 = -10$ milliers de dollars.
\end{solution}

\begin{solution}{exo:rc:rates-risk:2}
Elles décrivent deux scénarios différents. L'échéancier au pair additionne la variation due à une hausse
d'un point de base de chaque \emph{cotation} avec recalibrage, que les
swaps d'équilibrage ramènent à zéro~; l'échéancier zéro additionne la variation liée à une
hausse d'un point de base de chaque \emph{taux zéro de \omterm{def:rc:curve-construction:pillar}{pilier}}, qui déplace chaque
cotation au pair d'un peu plus d'un point de base (les taux diffèrent par la capitalisation et
la convention de décompte) et de montants différents, de sorte que les \omterm{def:rc:rates-risk:ladder}{sensibilités par pilier} du portefeuille
ne se compensent plus.
\end{solution}

\begin{solution}{exo:rc:rates-risk:3}
$100 - 98.2 = 1.8\%$ de la variance. L'écart-type quotidien du \omterm{def:rc:rates-risk:factors}{facteur de niveau}
est $\sqrt{\lambda_1} = 13.6$ points de base de score, soit un mouvement d'environ
$13.6\times0.40\approx5$ points de base au ventre de la courbe.
\end{solution}

\begin{solution}{exo:rc:rates-risk:4}
$367+713+1\,531+3\,205+5\,120+74\,016 = 84\,952$, contre un échéancier au pair de
$83\,778$~: l'échéancier zéro est supérieur dans le rapport $1.014$, facteur par lequel une
translation d'un point de base des taux zéro capitalisés en continu en ACT/365 fait varier
le taux au pair annuel ACT/360. Les \omterm{def:rc:curve-construction:pillar}{piliers} courts apparaissent parce que la hausse
d'un taux zéro précoce fait baisser la valeur des coupons fixes dus par le payeur~; dans l'espace des
cotations, le recalibrage absorbe cet effet dans le \omterm{def:rc:curve-construction:pillar}{pilier} à dix ans.
\end{solution}

\begin{solution}{exo:rc:rates-risk:5}
Les translations de taux clés sont définies sur la courbe zéro elle-même, quels que soient les instruments et
l'interpolation qui l'ont bâtie~; un échéancier au pair dépend de l'ensemble d'instruments et, hors
des \omterm{def:rc:curve-construction:pillar}{piliers}, de l'interpolation (chapitre~1). Il faut transmettre les taux clés (ou les
\omterm{def:rc:rates-risk:ladder}{sensibilités par pilier} propres à un régulateur) pour la comparabilité entre desks~; donner au trader l'échéancier au pair, dont les
\omterm{def:rc:rates-risk:ladder}{sensibilités par pilier} sont directement les notionnels des swaps au pair de couverture.
\end{solution}

\begin{solution}{exo:rc:rates-risk:6}
La couverture minimise la variance, et non les \omterm{def:rc:rates-risk:ladder}{sensibilités par pilier}~: les risques à sept ans et à dix ans ont
des signes opposés sur des échéances fortement corrélées, si bien qu'ensemble ils portent peu
de variance (le risque à dix ans compense à peu près celui à sept ans)~; les couvrir coûterait
du notionnel et de la fourchette pour presque aucune réduction du risque. Il s'agit d'une petite
position de spread sept ans contre dix ans.
\end{solution}

\begin{solution}{exo:rc:rates-risk:7}
Meilleur swap isolé~: celui à vingt ans, qui laisse 86.5~\% de la variance. Meilleur
ensemble de quatre~: deux, sept, vingt et trente ans, qui laissent 2.85~\%, contre 4.2~\%
pour les trois meilleurs~: le quatrième swap apporte peu.
\end{solution}

\begin{solution}{exo:rc:rates-risk:8}
Être neutre aux trois premières composantes laisse 68.9~\% de la variance de ce
portefeuille~: les trois facteurs décrivent comment toute la courbe se déplace, et non le
spread vingt ans contre trente ans où se loge le risque du portefeuille, que portent les
composantes supérieures. La neutralité aux facteurs n'élimine que le risque factoriel~; le résidu doit
être mesuré sur la matrice de covariance complète, ou couvert dans les \omterm{def:rc:rates-risk:ladder}{sensibilités par pilier} où il se trouve.
\end{solution}

\begin{solution}{pb:rc:rates-risk:1}
\textbf{1.} $g^\top\delta q = -\text{USD}~3\,000\,000$.
\textbf{2.} $-\text{USD}~3\,077\,970$~; l'écart, soit $-77\,970$, est du second ordre
(l'\cref{ex:rc:rates-risk:gamma}).
\textbf{3.} Trente ans, $+958\,308\times9 = +8.62$ millions, et vingt ans,
$-1\,135\,816\times7 = -7.95$ millions~; puis deux ans ($-2.76$ millions) et sept
($-2.00$).
\textbf{4.} Une pentification (taux courts en baisse, taux longs en hausse) avec un twist autour de dix ans~;
le portefeuille a perdu sur les \omterm{def:rc:rates-risk:ladder}{sensibilités par pilier} à deux et sept ans, qui y sont
longues en taux, et ses gains et pertes sur la partie longue de la courbe se sont presque compensés.
\textbf{5.} Il ne mesure que la somme des \omterm{def:rc:rates-risk:ladder}{sensibilités par pilier}, c'est-à-dire la réponse à un mouvement que personne
n'a vu ce jour-là.
\textbf{6.} $+21.2$.
\textbf{7.} $-\text{USD}~3.09$ millions.
\textbf{8.} $21.2/4.87 = 4.3$ écarts-types~: une journée rare de mouvement de pente.
\textbf{9.} Mouvements d'échéances isolées par rapport à leurs voisines~: les
grandes \omterm{def:rc:rates-risk:ladder}{sensibilités par pilier} opposées à vingt et trente ans les rendent importants.
\textbf{10.} Approximativement pour la partie pente~; pas pour le reste, puisque le risque du
portefeuille hors des trois facteurs est important (exercice~8).
\textbf{11.} USD~1\,288\,843 par jour.
\textbf{12.} 58.3~\%.
\textbf{13.} Recevoir le fixe sur USD~1.196 milliard à deux ans, payer le fixe sur USD~801
millions à vingt ans, recevoir le fixe sur USD~515 millions à trente ans.
\textbf{14.} USD~264\,554 par jour ($\sqrt{0.042}\times1\,288\,843$).
\textbf{15.} $-\text{USD}~315\,044$ (premier ordre)~: la couverture a supprimé les neuf dixièmes
de la perte.
\textbf{16.} Chaque couverture coûte de la fourchette et du bilan, et annuler chaque \omterm{def:rc:rates-risk:ladder}{sensibilité par pilier}
supprime des compensations que les mouvements historiques rendent presque sans risque~; un critère
de variance couvre ce qui bouge réellement.
\textbf{17.} Une fois par mois ou par semaine, sur une fenêtre assez longue pour des estimations stables
(plusieurs années), mais pondérée vers les données récentes, et contrôlée selon le
régime (politique en pause, en baisse ou en hausse)~; soumettre la couverture à des tests de résistance sur d'autres fenêtres.
\textbf{18.} L'échéancier, les expositions aux facteurs, la variance résiduelle après les
couvertures en place, le P\&L sur un ensemble de journées historiques de twist, et le gamma.
\textbf{19.} Résultat nommé~: \emph{la journée de twist} coûte au portefeuille neutre en DV01 USD~3.00
millions au premier ordre (3.08 par réévaluation complète)~; des swaps recevant le fixe à deux ans, payant
le fixe à vingt ans et recevant le fixe à trente ans laissent 4.2~\% de sa variance quotidienne.
\textbf{20.} C'est une direction parmi beaucoup d'autres, et la courbe ne bouge presque jamais
exactement dans cette direction.
\end{solution}

\begin{solution}{iq:rc:rates-risk:1}
Seulement contre les mouvements parallèles. Il faut regarder l'échéancier par \omterm{def:rc:curve-construction:pillar}{pilier}~: une somme nulle peut masquer
de grandes \omterm{def:rc:rates-risk:ladder}{sensibilités par pilier} qui se compensent, un steepener ou un butterfly, qui perdent sur les
mouvements qui se produisent réellement. Il faut demander l'échéancier, les expositions aux facteurs et un ensemble de
scénarios historiques de twist.
\iqlookfor{la différence entre un risque scalaire et un risque vectoriel.}
\end{solution}

\begin{solution}{iq:rc:rates-risk:2}
Avec le jacobien $J = \partial(\text{cotations du modèle})/\partial(\text{zéros aux
piliers})$ issu de la calibration~: $g_z = J^\top g_q$, et $g_q =
(J^\top)^{-1}g_z$. Ou par force brute~: décaler les cotations et recalibrer, ou décaler
les \omterm{def:rc:curve-construction:pillar}{piliers} sans recalibrer. Il faut vérifier que les deux voies concordent.
\iqlookfor{la règle de dérivation en chaîne à travers la calibration.}
\end{solution}

\begin{solution}{iq:rc:rates-risk:3}
Le niveau (toutes les échéances de même signe), la pente (court contre long) et la courbure
(corps contre ailes)~; sur les variations quotidiennes des bons du Trésor américain entre 2016 et 2026, elles expliquent
environ 85~\%, 11~\% et 2~\%, soit 98~\% ensemble. Ces parts dépendent du marché et
de la période (la partie courte de la courbe se comporte différemment autour des décisions de politique monétaire).
\iqlookfor{les formes et les ordres de grandeur, et le fait qu'il s'agit d'estimations.}
\end{solution}

\begin{solution}{iq:rc:rates-risk:4}
Calculer l'échéancier du portefeuille et la covariance des variations de cotations~; pour chaque
triplet candidat de swaps liquides, résoudre la couverture de variance minimale et retenir le
triplet qui laisse la plus faible variance résiduelle, en vérifiant qu'il tient sur une autre
fenêtre. Privilégier les \omterm{def:rc:curve-construction:pillar}{piliers} liquides, pénaliser les notionnels importants. Ici, dans le manuel,
deux, dix et trente ans laissent 58~\%, et deux, vingt et trente ans en laissent 4~\%.
\iqlookfor{un critère, non une règle empirique.}
\end{solution}

\begin{solution}{iq:rc:rates-risk:5}
Termes du second ordre~: la matrice des gammas, y compris les \omterm{def:rc:rates-risk:crossgamma}{gammas croisés} entre \omterm{def:rc:curve-construction:pillar}{piliers},
et, pour les portefeuilles d'options, le vanna et le volga. Les calculer par différences centrées de
l'échéancier ou de la valeur, et ajouter $\tfrac12\delta q^\top\Gamma\delta q$~; s'il subsiste un
résidu, chercher les facteurs de risque manquants (base, volatilité) ou des
données d'entrée périmées.
\iqlookfor{le développement de Taylor à l'ordre deux, puis les facteurs manquants.}
\end{solution}

\begin{solution}{iq:rc:rates-risk:6}
Lorsque le risque du portefeuille se situe en dehors des premiers facteurs (positions de spread entre
échéances voisines)~; lorsque le régime de la fenêtre d'estimation diffère de celui d'aujourd'hui
(taux nuls, puis hausses)~; lorsque les facteurs sont instables en signe ou en
ordre~; et, les jours de stress, lorsque les corrélations changent. Toujours rapporter le résidu
sur la matrice de covariance complète.
\iqlookfor{l'ACP comme modèle, avec son propre risque.}
\end{solution}
